% =====================================================================================================
\section{Supplement: further results}\label{sec:supp}

Each result below comes with one line on why it matters; the status notes give their standing.

\subsection{The posterior}\label{sec:supp-post}

\why{Without the beta-min condition, the posterior still rejects over-fitted models and still keeps every true predictor
that is strong enough; specificity, however, can fail.}

\begin{remark}[Without beta-min]\label{rem:no-betamin}
Under the other assumptions of Theorem~\ref{thm:post-fixed}, the posterior mass of models that strictly contain $\Sz$ is
$\OP(\rn)$, and every true predictor with $\theta_j^{*2}\ge c\,\sigs^2\log n/(n\varphi_0^2)$, $c>c_0(\sz,A)$, has
$1-\PIP_j=\OP(\rn)$. Specificity can fail without beta-min: a numerical counterexample, reported with the proof, shows this.
\end{remark}

\status{Part of the posterior concentration result, with the status of the in-probability form of
Theorem~\ref{thm:post-fixed}.}

\why{It replaces the eigenvalue floor away from the true model, and the beta-min condition, by a floor on the true model
alone and a margin on the unexplained signal; it keeps the vanishing floor \as{NC}.}

\begin{localassumption}[Floor on the true model and an identifiability margin]\label{asm:margin}
\begin{asmitems}
  \item $\lmin(\Rhat_{\Sz})\ge\varphi_0^2$ for a fixed $\varphi_0\in(0,1]$.
  \item Every model $S\in\Om$ that misses a true predictor and has at most $j_0:=\lceil(2A+1)\sz\rceil+1$ predictors
    outside $\Sz$ leaves unexplained signal $\norm{(I-P_S)\mu}^2\ge c_{\mathrm{id}}\,\sigs^2\log n$, for a fixed
    $c_{\mathrm{id}}>c_1(\sz,A)$, an explicit constant (Appendix~\ref{app:constants}).
\end{asmitems}
\end{localassumption}

\begin{theorem}[Posterior concentration under an identifiability margin]\label{thm:post-margin}
Let the design be fixed and assume \as{GL}, \as{DIM}, \as{SP}, \as{BC}, \as{NC} and Assumption~\ref{asm:margin}. Then
the conclusions~\eqref{eq:post-op}--\eqref{eq:post-exp} of Theorem~\ref{thm:post-fixed} hold, the expectation bound with
constants depending only on $(\sz,A,\sigs^2,\bar\theta,\varphi_0,g_0,c_{\mathrm{id}})$.
\end{theorem}

\status{The in-probability conclusions: \textbf{proved and independently checked}. The expectation bound:
\textbf{proved and checked, conditional on the same supporting result} (the lower-bound proposition, not yet independently
checked) as the expectation bound of Theorem~\ref{thm:post-fixed}. Registry: Appendix~\ref{app:status}, row 1.1.}

Under \as{EF} and \as{BM}(a), Assumption~\ref{asm:margin} holds with $c_{\mathrm{id}}=C_\beta^2\varphi_0^2(1-1/n)$. So
Theorem~\ref{thm:post-margin} contains Theorem~\ref{thm:post-fixed} when $C_\beta^2\varphi_0^2>c_1$; when
$c_0<C_\beta^2\varphi_0^2\le c_1$, Theorem~\ref{thm:post-fixed} applies and Theorem~\ref{thm:post-margin} is not known to
apply: the two coexist, and neither contains the other.

\why{It says what the posterior selects when some true predictors are missing from the data: the predictors used by the
best linear predictor built from the available ones.}

In this result the method is run on the response $\yr$ and a set $\mathcal C$ of $p$ candidate predictors, with rows
$x_{\mathcal C,i}$ collected in $x_{\mathcal C}$; $\Sigma$ and $R$ are the covariance and the correlation matrix of a
candidate row.

\begin{localassumption}[Candidate-only fitting]\label{asm:omitted}
\begin{asmitems}
  \item \emph{Model.} The candidate rows are i.i.d.\ $\Normal_p(m_{\mathcal C},\Sigma)$, and, given them, the responses are
    independent with $\yr_i\mid x_{\mathcal C}\sim\Normal(\alpha'+x_{\mathcal C,i}^\top\beta^{\mathcal C},\sigma^2_{\mathcal C})$,
    $\sigma^2_{\mathcal C}>0$. We write $\Sz^{\mathcal C}:=\supp\beta^{\mathcal C}$ and $s_{\mathcal C}:=\abs{\Sz^{\mathcal C}}$.
  \item \emph{Floor.} $\lmin(R_T)\ge c_{\mathcal C}$ for every $T\subseteq\mathcal C$ with
    $\abs{T}\le3(A+1)(s_{\mathcal C}+1)\log n$, for a fixed $c_{\mathcal C}>0$.
  \item \emph{Signal relative to the residual variance.} $\beta^{\mathcal C\top}\Sigma\beta^{\mathcal C}\le
    M_{\mathcal C}\sigma^2_{\mathcal C}$ for a fixed $M_{\mathcal C}$, and $n\min_{j\in\Sz^{\mathcal C}}\Sigma_{jj}(\beta^{\mathcal
    C}_j)^2/(\sigma^2_{\mathcal C}\log n)\to\infty$.
\end{asmitems}
\end{localassumption}

\begin{theorem}[Omitted variables]\label{thm:omitted}
Assume \as{DIM} (with $p$ the number of candidates), Assumption~\ref{asm:omitted}, and that $s_{\mathcal C}\ge1$ is fixed. No
condition is placed on how $m_{\mathcal C}$, $\Sigma$, $\beta^{\mathcal C}$ or $\sigma^2_{\mathcal C}$ vary with $n$. Then
\[
  \E\big[1-\pi(\Sz^{\mathcal C}\mid y)\big]=O\Big(\sqrt{\frac{\log n}{n}}\Big),\qquad
  \P\big(\HPM\neq\Sz^{\mathcal C}\big)=O\Big(\sqrt{\frac{\log n}{n}}\Big).
\]
\end{theorem}

\status{\textbf{Proved and checked, conditional on} the expectation bound of Theorem~\ref{thm:post-fixed} (hence on the
lower-bound proposition, not yet independently checked) and on two lemmas on Gaussian designs that are independently
verified. Registry: Appendix~\ref{app:status}, row 1.4.}

For example, if candidates and omitted predictors are jointly Gaussian, $\beta^{\mathcal C}=\beta^*_{\mathcal C}+
\Sigma_{\mathcal C\mathcal C}^{-1}\Sigma_{\mathcal C\mathcal M}\beta^*_{\mathcal M}$: each candidate absorbs the part of the
omitted signal it can predict, and the rest is added to the noise variance $\sigma^2_{\mathcal C}$.

\subsection{The distribution}\label{sec:supp-dist}

\why{It shows what the collinearity condition of Result~2 buys: dropping it costs one logarithmic factor in the radius,
and the posterior also puts almost no mass on large models.}

\begin{proposition}[Level-$M$ form, and the case without the collinearity condition]\label{prop:dist}
\textup{(i)} There is a constant $C_{\mathrm s}$, depending only on $(s,A,L,c_\psi,\bar b)$, such that the following holds. The
constants of Theorem~\ref{thm:dist} serve every level $M$ with $M_0\le M\le\sqrt n/\log n$, and $M$ may depend on $n$:
whenever $n\ge n_0$ and Assumption~\ref{asm:dist} holds with the window $q_n$ of this $M$,~\eqref{eq:dist-i} holds, and also
\begin{equation}\label{eq:dist-ii}
  \P^*\Big[\,\Pi\big(\abs{S}>s+C_{\mathrm s}M\ \big|\ \Dn\big)>n^{-c_{\mathrm H} M}\Big]\ \le\ n^{-c_{\mathrm H} M}.
\end{equation}
The two exceptional sets can be taken to be the same, and the bounds are uniform over all laws of the data that satisfy
Assumption~\ref{asm:dist} with the given constants.
\textup{(ii)} If Assumption~\ref{asm:dist}\ref{asm:dist-nec} is dropped and Assumption~\ref{asm:dist}\ref{asm:dist-subg} is
required only for single predictors, then~\eqref{eq:dist-i} holds with radius $\sqrt{M\log n\cdot\log(\praw\vee n)/n}$,
and~\eqref{eq:dist-ii} with $s+C_{\mathrm s}M\log n$ in place of $s+C_{\mathrm s}M$.
\end{proposition}

\status{\textbf{Proved and independently checked} (part of the same result as Theorem~\ref{thm:dist}). Registry:
Appendix~\ref{app:status}, row 2.1.}

\why{With a fixed eigenvalue floor and every true signal above the detection scale, the distribution is learned at the
parametric rate.}

\begin{theorem}[Parametric rate when the true predictors are detectable]\label{thm:dist-param}
Assume Assumption~\ref{asm:dist} with the window $q_n$ of level $M=1$, with exactly $s\ge1$ nonzero coordinates of $\bstar$;
that $\lmin(\Corr(x_T))\ge\psi$ for a fixed $\psi>0$ and every set $T$ of at most $(2A+3)s+2$ predictors; and that
$\min\{\abs{\beta^*_j}\sd(x_j)/\sigs:\beta^*_j\neq0\}\ge b_n$, where $b_n>0$ and $b_n^2n/\log n\to\infty$. Then
$\Pi(d_H(f,f^*)>M_n/\sqrt n\mid\Dn)\to0$ in $\P^*$-probability, for every sequence $M_n\to\infty$.
\end{theorem}

\status{\textbf{Proved and independently checked}. Registry: Appendix~\ref{app:status}, row 2.1.}

\why{The distribution guarantee carries over to the list of models that a completed run returns, whenever that list
holds a non-negligible share of the posterior.}

A completed run returns a list $\mathcal T$ of at most $40p+1$ distinct models with their log scores; let
$\pi(\mathcal T\mid\Dn)$ be the posterior probability of the list, and $\hat\Pi$ the posterior restricted to it (the models
of $\mathcal T$ weighted in proportion to their exponentiated log scores, each with its exact conditional posterior of
$(\beta,\sigma^2)$; $\hat\Pi$ is not computed by the method).

\begin{corollary}[The output of the sampler]\label{cor:dist-output}
Assume the conditions of Proposition~\ref{prop:dist}(i), and that the returned log scores are computed in exact
arithmetic. With $\P^*$-probability at least $1-n^{-c_{\mathrm H}M}$, every completed run, whatever its random numbers, with
$\pi(\mathcal T\mid\Dn)\ge n^{-c_{\mathrm H}M/2}$ satisfies $\hat\Pi(d_H(f,f^*)>\sqrt{M\log n/n})\le n^{-c_{\mathrm H}M/2}$. If a
completed run scored $\Sz$ and $\pi(\Sz\mid\Dn)>\tfrac12$, then the list contains $\Sz$, $\Shat=\Sz$, and $\pi(\mathcal
T\mid\Dn)>\tfrac12$.
\end{corollary}

\status{\textbf{Proved and checked, conditional on} exact arithmetic and on \samplerconds; the description of the
output was checked against the code. Registry: Appendix~\ref{app:status}, row 2.1.}

Whether the list carries enough posterior probability is a property of the run; the corollary does not assert it.

\subsection{The method}\label{sec:supp-method}

\why{It explains why the sampler starts in the right place: without leakage, the start law already concentrates on the
true model.}

\begin{theorem}[The start law concentrates]\label{thm:start}
Assume \as{GR}, \as{DIM} with $p>\sz$, \as{FSt}, \as{V}, \as{NL}, and $\lmin(\Sigma_{n,S})\ge\varphi_{\mathrm b}^2$ for every
$S$ with $\abs{S}\le16+\sz$, for a fixed $\varphi_{\mathrm b}>0$. Then $\P(\pi_0(\Sz)\ge1-n^{-c})\to1$ for every
$c<\tfrac12$. If moreover $\Cov(W_l,W_j)=0$ for every null $l$ and true $j$, then $1-\pi_0(\Sz)=\OP(\rn)$. The exponent
$\tfrac12$ cannot be improved: $1-\pi_0(\Sz)\ge\tfrac13(n+1)^{-1/2}$ for every valid data set with $n\ge4$ and $\sz<K$.
\end{theorem}

\status{\textbf{Proved and independently checked}. Registry: Appendix~\ref{app:status}, row 3.3.}

With leakage, the start law concentrates instead on the models between the true model and the leakage set: keep the
other assumptions of Theorem~\ref{thm:start}, and let \as{NL} hold with $\Sz$ replaced by a deterministic $Q\supseteq\Sz$ with $\abs{Q}\le K$, so
that only predictors outside $Q$ are required not to leak. Then $\P\big(\pi_0(\Lat(Q))\ge1-(\abs{Q}+1)n^{-1/2}e^{C\sqrt{\log n}}\big)\to1$.

\why{It is the only reported-model statement that does not assume Gaussian rows: its conditions are on the sample
$(X,\tstar)$.}

For a set $U$ and a predictor $l\notin U$, the \emph{signal gain} $\delta^2_{U,l}:=\norm{(P_{U\cup\{l\}}-P_U)\mu}^2$ is the
increase in explained mean signal when $x_l$ joins $X_U$; the \emph{residual scale} is $\tau^2_U:=\sigs^2+\norm{(I-P_U)\mu}^2/(n-1)$;
and $\hat\beta_{U,l}:=(X_U^\top X_U)^{-1}X_U^\top x_l$, padded with zeros.

\begin{localassumption}[Fixed-design visibility and leakage]\label{asm:reported-fixed}
\begin{asmitems}
  \item\label{asm:rf-floor} \emph{Floors.} $\lmin(\Rhat_T)\ge\varphi_0^2$ for every $T$ with $\abs{T}\le16$ and for every
    $T=Q\cup\{l\}$, $l\in\set{p}$, with $Q$ as in~\ref{asm:rf-leak}.
  \item\label{asm:rf-vis} \emph{Visibility.} There is a deterministic $\omega_n\to\infty$ such that
    $\delta^2_{U,j}\ge\omega_n\log n$ for every $j\in\Sz$ and every $U\not\ni j$ with $\abs{U}\le15$.
  \item\label{asm:rf-leak} \emph{Bounded leakage.} There are a deterministic $Q\supseteq\Sz$ (with $\abs{Q}\le K$ and
    ($Q=\set{p}$ or $K\ge2\abs{Q}+32\sz+1$) for large $n$) and a fixed $\eta\in[0,\eta_A)$,
    $\eta_A:=(\sqrt{2A+1}-\sqrt{2A})^2$, such that $\delta^2_{U,l}\le\eta\,\tau^2_U\log n$ for every $l\notin Q$ and every
    $U\not\ni l$ with $\abs{U}\le15$.
  \item\label{asm:rf-size} \emph{Size of the leakage set.} One of: (d1) $\abs{Q\setminus\Sz}\le(1-\delta_Q)\log n$ for a
    fixed $\delta_Q\in(0,1)$; or (d2) $\abs{Q}\le n^b$ with $b<\tfrac12-\tfrac\eta2-\sqrt{2A\eta}$,
    $\theta^{*2}_{\min}n/(\abs{Q}+\log n)\to\infty$, and local regressions: for every $l$ a deterministic set
    $N_l\subseteq Q\setminus\{l\}$ with $\abs{N_l}\le h\log_2n$, for a fixed $h$ with $h+b<\tfrac12$, and a deterministic $L_n\to0$ with
    $L_n^2\log(\abs{Q}+1)=o(\log n)$, such that $\norm{\hat\beta_{T,l}-\hat\beta_{T\cap N_l,l}}_1\le L_n$ for every
    $T\in\Lat(Q)$ with $\abs{T}<K$ and $l\notin T$.
\end{asmitems}
\end{localassumption}

\begin{theorem}[The reported model equals the true model: fixed design]\label{thm:reported-fixed}
Let the design be fixed, with no raw predictor constant in the sample (an input condition of the delivery result; then
$p=\praw$). Assume \as{GL}, \as{DIM}, \as{SP}, \as{BC},
\as{EF}, \as{NC}, \as{BM}(b), Assumption~\ref{asm:reported-fixed} and Convention~\as{C1}. Then the conclusions of
Theorem~\ref{thm:reported} hold: there are data events $\mathcal G_n$ with $\P(\mathcal G_n)\to1$ on which every
admissible run that produces output reports $\Shat=\Sz$, for every seed and calibration outcome; and the delivery
clauses hold under \as{C2} on $\Dom$ and under \as{C2p} for every admissible $n\ge3$.
\end{theorem}

\status{\textbf{Proved and independently checked}, with no open condition: the version proved in the registry is a stronger form of
this statement. Registry: Appendix~\ref{app:status}, row 3.2.}

\why{A certified small model rarely misses a true predictor, under much weaker conditions than Result~3; it is the
guarantee about the coded method that remains on pervasive-factor designs, which Result~3 excludes.}

\begin{theorem}[Certified screening]\label{thm:screen}
Let the design be fixed and assume \as{GL} and Convention~\as{C1}; \as{SP} is not assumed. Let $k_n\in\{0,\dots,K-1\}$ be
a model-size budget, and assume $\lmin(\Rhat_S)\ge\kappa$, for a fixed $\kappa>0$, for every $S$ with $\abs{S}\le k_n+\sz$, $\tmin\ge\tlow>0$ and
$\norm{\tstar}_1\le M_1$, with $\kappa$, $\tlow$, $M_1$, $\sigs^2$ fixed, and $(k_n+1)\log p+\log n\le c_{\mathrm{SB}}\,n$, where
$c_{\mathrm{SB}}:=\kappa^2\tlow^2/\{200(M_1^2+\sigs^2)\}$. Then
\[
  \P\big(\text{the run produces output},\ \Shat\text{ is certified},\ \abs{\Shat}\le k_n,\ \Sz\not\subseteq\Shat\big)\to0 .
\]
\end{theorem}

\status{\textbf{Proved and independently checked}. Registry: Appendix~\ref{app:status}, row 3.7.}

No visibility, leakage, start or block condition is used. A second form replaces the signal bounds by a margin on the
noiseless gain of a missing true predictor, with no condition on $p$, $\sz$ or the null predictors (also proved and
independently checked). The bound is on an intersection, not conditional on certification, and a run is not claimed to certify a small model
with probability tending to one. In existing numerics (not a theorem) every completed run certified its reported model:
2,394 of 2,394 in the 2,400-fit study, 376 of 376 in the strong-factor study, 71 of 71 in a screening study. On the
strong-factor designs with equicorrelation $\varrho\in\{0.5,0.7\}$, $p\ge2.5n$ and $\sz\in\{4,8\}$, the certified model contained every true predictor in 61 of 122 fits at $n=50$, 18 of 20 at $n=75$
and 30 of 30 at $n=100$; at the constants of these designs the size budget of Theorem~\ref{thm:screen} requires
$n\gtrsim1{,}450$.

\why{It is the only recorded route to the reported inclusion probabilities and median model that does not need \as{NL};
its global-overlap premise is not yet derived from conditions on the design.}

\begin{remark}[An alternative route for the estimates]\label{rem:alt-route}
An argument that does not go through the evidence estimate gives the conclusions of Theorem~\ref{thm:outputs} for the
inclusion probabilities and the median and highest posterior models under a global-overlap condition (stated in full with
the proof).
\end{remark}

\status{\textbf{Proved and independently checked}; its global-overlap premise is a named open premise. Registry:
Appendix~\ref{app:status}, row 3.6.}


\subsection{The calibration and the Markov moves in full}\label{sec:supp-cal}

\why{Every result holds for every outcome of the calibration, so the main text uses only three facts about it; this is the
procedure itself, for reproducibility.}

At tier $i$ the calibration's islands have $\Nc_i:=4^i\Nc_0$ particles, and the \emph{step ceiling} is
$\Tmax_i:=4^i\max(1000,\ell_n^2)$.

\paragraph{Stores of models.} Throughout the run the method keeps three stores of the models scored so far: the $40p+1$
best by $F$; the $20p$ best by $\gamma_0$; and the $20p$ best by $\gamma_1=e^{F}/m(\emptyset)$, that is, by $F$. The two
stores of $20p$ feed the candidate set below. \emph{The best model found so far} is the model with the highest $F$ among
the models recorded so far: during the calibration, the block-table models (recorded before the pilot starts) and those
scored in the calibration; in the final run, the block-table models, those scored in the final run, and the models passed
on by the calibration.

\begin{center}
\fbox{\begin{minipage}{0.93\linewidth}\small
\textbf{Algorithm 1: calibration.}
\begin{enumerate}[leftmargin=1.6em,itemsep=1pt,label=\arabic*.]
  \item Set the tier $i\leftarrow0$.
  \item \emph{Pilot.} Run one island of $\Nc_i$ particles (reweight, resample, move, as in Section~\ref{sec:algorithm})
    with an adaptive schedule. From the current $\xi$, the next temperature is $\xi'=1$ if that step keeps an effective
    sample size of at least $0.90\,\Nc_i$; otherwise it is found by 60 bisections on $(\xi,1]$. If the step is shorter than
    $1/\Tmax_i$, take $\xi'=\min(1,\xi+1/\Tmax_i)$ instead and count a \emph{forced step}. At $\xi=1$, run $\Lbi$ burn-in rounds and
    four measured rounds, then a local search (move cap $\ell_n$) from the best model found so far.
  \item \emph{Grid.} Build a grid from the pilot's temperatures on the candidate set $\mathcal V$ (below). If the
    construction fails, record a trigger and skip step~4.
  \item \emph{Rebuild loop.} Repeat at most $\ell_n$ times: run four diagnostic islands of $\Nc_i$ particles on the
    current grid (each as in Section~\ref{sec:algorithm}, with $\Lbi$ burn-in and four measured rounds); run a local search
    (move cap $\ell_n$) from the best model found so far; recompute $\mathcal V$, and rebuild the grid from the pilot's
    temperatures. Stop when the rebuilt grid is identical to the previous one (the grid has \emph{converged}), or when a
    rebuild fails.
  \item \emph{Triggers.} If $i=0$ and any trigger fires (list below): repeat the check before the run, set $i\leftarrow1$,
    and go to step~2. The scored models, the three stores and the score cache carry over; the tier-$0$ temperatures do
    not, and tier~1 uses its own random streams. At tier~1, triggers are only recorded.
  \item \emph{Fallback.} If the grid construction failed at tier $1$, use the last grid built successfully at tier $1$,
    or else the tier-$1$ pilot's own temperatures. (An unconverged grid is used as it is.)
  \item \emph{Freeze} the grid, the tier $i$, the number of particles $N_i$ and $\Lbi$, and pass the three stores and the
    calibration's best model to the final run, as candidates for the reported model (they are not particles). With the
    fallback on (the default), a grid always exists at this point; if it does not, stop the run without estimates.
\end{enumerate}
\end{minipage}}
\end{center}

\paragraph{The candidate set and the grid.} The candidate set $\mathcal V$ is the empty model together with the two stores
of $20p$ models (the best by $\gamma_0$ and the best by $F$). On $\mathcal V$, let
$\psi(u):=\log\sum_{S\in\mathcal V}\gamma_0(S)e^{u\,\Delta(S)}$. A step from temperature $u$ to $v$ has \emph{restricted effective
sample size}
\[
  \rho(u,v):=\exp\{2\psi(v)-\psi(u)-\psi(2v-u)\},
\]
the effective-sample-size fraction that the step would leave if the particles were the models of $\mathcal V$ with
weights proportional to $\gamma_u$. The grid \emph{refines} the pilot's temperatures. Starting at $u=0$, let $b$ be the
smallest pilot temperature above $u$. If $\rho(u,b)\ge0.90$, the next grid point is $b$; otherwise it is the largest
point of $(u,b]$ that passes, found by 60 bisections. The construction continues from the new point, again aiming at the
smallest pilot temperature above it, until it reaches $1$; so every pilot temperature is in the grid, and points are
inserted where a pilot step fails the check on $\mathcal V$. The construction \emph{fails} if a step
makes no numerical progress, if an appended temperature does not pass the $0.90$ check (temperatures within $10^{-12}$ of
$1$ are set to $1$), or if more than $\Tmax_i$ steps are needed.

\paragraph{Triggers.} At tier $0$ the calibration moves to tier $1$ if any of these occurs: \emph{island disagreement}
(E1: some diagnostic island's best log score is more than $1$ below the best score found so far); \emph{evidence
spread} (E2: the sample standard deviation of the four islands' log-evidence estimates exceeds $1$); \emph{late
improvement} (E3: a calibration local search, the pilot's included, ended with a strictly higher best log score than the
best before it started); a forced pilot step; a failed grid construction; a grid
that did not converge; or a local search that was still improving when it reached its move cap. E1 and E2 are taken from
the last completed rebuild; E3 and the move-cap trigger count over the pilot and every rebuild.

\paragraph{When the calibration succeeds.} The calibration \emph{succeeds} when it freezes a grid, built by the main
construction or by the fallback. Without a trigger, tier $0$ freezes its own grid. The calibration's own estimates are
archived and never reported.


\paragraph{The Markov moves.} One \emph{round} applies $p+B+2$ moves to a particle, in an order drawn uniformly at random
afresh for each particle and each round; each move leaves $\pi_\xi$ invariant.
\begin{itemize}[leftmargin=1.5em,itemsep=1pt]
  \item \emph{Inclusion update of predictor $j$} (one for each $j\in\set{p}$; a Gibbs update): with
    $a:=\gamma_\xi(S\cup\{j\})/\gamma_\xi(S\setminus\{j\})$, include $j$ with probability $o_j:=a/(1+a)$, or $o_j:=0$ if
    including $j$ would exceed the cap $K$. The probability $o_j$, not the drawn indicator, is what the inclusion-probability
    estimate uses (step~3 below).
  \item \emph{Uniform swap:} pick a member $j$ of the model and a non-member $l$ uniformly at random; propose
    $S':=(S\setminus\{j\})\cup\{l\}$, and accept with probability $\min\{1,\gamma_\xi(S')/\gamma_\xi(S)\}$. (If
    the model is empty or equal to $\set{p}$, the move does nothing.)
  \item \emph{Neighbour swaps} ($B$ of them): pick a member $j$ uniformly, and a non-member $l$ uniformly among $j$'s
    neighbours in the neighbour graph; propose the swap, and accept with the Metropolis--Hastings probability, which
    includes the correction $\abs{\mathcal N(j)\setminus S}/\abs{\mathcal N(l)\setminus S'}$ for the different
    neighbourhood sizes. (If the model is empty, or $j$ has no neighbour outside it, the move does nothing.)
  \item \emph{Independence move:} propose a fresh draw $S'$ from $\pi_0$, and accept with probability
    $\min\{1,\exp\{\xi\,(\Delta(S')-\Delta(S))\}\}$.
\end{itemize}
Swaps keep the model size; the inclusion updates and the independence move can change it. Every move stays in $\Om$:
the inclusion updates enforce the cap, and $\pi_0$ lives on $\Om$.


\subsection{The engine for \texorpdfstring{$p=10^6$}{p=10\^{}6} (an engineered variant)}\label{sec:supp-1m}

\why{It is the engine of the planned runs at $p=10^6$; it differs from the method, and no theorem covers it.}

\emph{This describes a different engine, not the method of Section~\ref{sec:method}, which the results are about.} To run at $p=10^6$, the planned engine changes the method as listed in Table~\ref{tab:1m}, taken from the
draft engineering specification of 2026-11-15; entries marked ``to be decided'' depend on a precision study. The run at
$p=10^6$ will be reported as an engineered variant validated against the method above on shared problem sizes ($p=10^4$,
and $p=10^5$ where possible), not as a run covered by the theorems.

{\small
\begin{longtable}{@{}>{\raggedright\arraybackslash}p{0.2\linewidth}>{\raggedright\arraybackslash}p{0.33\linewidth}>{\raggedright\arraybackslash}p{0.4\linewidth}@{}}
\caption{The $10^6$ engine against the method.}\label{tab:1m}\\
\toprule
Aspect & The method & The $10^6$ engine \\ \midrule \endhead
particles per island & $4^i\max(2500,n,p)$ & fixed, $2{,}500$ ($10^4$ in one arm of the precision study) \\
islands & four, in one process & sixteen, in separate processes, merged afterwards (union of stored models; equal-weight inclusion probabilities as the primary estimate) \\
random numbers & sequential streams & addressed by position (counter-based) \\
scoring & the score of Section~\ref{sec:posterior} & a batched kernel with different floating-point arithmetic, agreeing to about $10^{-9}$ in relative terms \\
blocks & the partition of Section~\ref{sec:algorithm}, $B\le16$ & complete linkage truncated at a correlation threshold, the same top-down cut, packing by smallest index; $B\le14$ \\
calibration & Algorithm~1, with at most one move to tier 1 & the pilot and at most a fixed number of rebuilds (0 or 3, to be decided), no move to tier 1 except after a failed construction \\
grid & as frozen & optionally each interval split into $m\in\{1,4\}$ equal parts (to be decided) \\
estimates & those of Section~\ref{sec:outputs-method} & the same, plus a deterministic lower bound on $\log Z$ from the stored models, an evidence estimate based on how often the islands visit the stored models, and inclusion probabilities computed from the stored models \\
resampling and moves & multinomial; $B$ neighbour swaps per round & systematic resampling and more neighbour swaps per round, off unless the precision study shows a gain \\
caches & a score cache, the store of the $40p+1$ best models, and the two stores of $20p$ & no score cache and no stores of $20p$ in the final islands; each island keeps a store of its $40p+1$ best models \\
\bottomrule
\end{longtable}}


\subsection{Open questions}\label{sec:open}

None of these is needed for any statement above.
\begin{enumerate}[leftmargin=1.5em,itemsep=1pt]
  \item Remove the vanishing floor \as{NC} from Theorem~\ref{thm:post-margin}, and weaken its floor on the true model.
  \item Prove $\E[1-\pis{\Sz}]\gtrsim\rn$ on every design with a floor on the sets $\Sz\cup\{j\}$.
  \item Whether Assumption~\ref{asm:dist}\ref{asm:dist-nec} is needed for the rate $\rn$ of Result~2.
  \item Leakage sets of size $n^b$ with $b\ge\tfrac12$, and raw-scale signals in Results~3 and~4.
  \item Consistency of the reported estimates on designs with leakage (AR and block designs).
  \item The parallel method: an inheritance statement for an output-identical engine, and results for many islands and
    island weighting at very large $p$.
\end{enumerate}
